\documentclass[11pt,a4paper]{article}
\usepackage[T1]{fontenc}
\usepackage{lmodern,amsmath,amssymb,amsthm,mathtools,microtype,needspace}
\usepackage[margin=25mm]{geometry}
\usepackage[colorlinks=true,allcolors=blue]{hyperref}
\newtheorem{theorem}{Theorem}
\newtheorem{lemma}{Lemma}
\newcommand{\N}{\mathbb N}
\newcommand{\itf}[2]{F_{#1}^{\,#2}}
\newcommand{\itl}[2]{L_{#1}^{\,#2}}
\title{Counterexamples to a late-iterate ordering conjecture for generalized Hofstadter functions}
\author{Dev D. Goyal}
\date{26 October 2026}
\usepackage[numbers]{natbib}
\usepackage{xurl}
\begin{document}
\maketitle
\begin{abstract}
Letouzey, Li, and Steiner conjecture a pointwise order between iterates
of successive generalized Hofstadter functions. This paper disproves that
order. A finite witness gives $F_4^9(132)=8>7=F_5^{10}(132)$.
An inverse recurrence gives a counterexample for every parameter $k\geq4$,
at the first iteration index $j=2k+1$ in the conjectured range.
The proof uses an exact difference of $k-3$ between two inverse iterates.
The Lean formalization constructs the functions on all natural numbers
from their defining recurrence. It proves both the finite witness and
the full parameter family.
\end{abstract}
\section{The conjecture and a finite counterexample}

For \(k\ge1\), define
\[
 F_k(0)=0,\qquad F_k(n+1)=n+1-\itf{k}{k}(n).
 \tag{1}
\]
Here \(f^r\) means \(r\)-fold composition, with \(f^0\) the identity.
These are the functions in Section 1.1 of \cite{LLS}.
The unnumbered conjecture immediately before Lemma 7.6 in arXiv version 2
asserts
\[
 \forall k\ge1\ \forall j>2k\ \forall n\in\N,
 \qquad \itf{k}{j}(n)\le\itf{k+1}{j+1}(n).
 \tag{C}
\]
One counterexample is sufficient to disprove this universal statement.

\begin{theorem}[Finite counterexample]
The conjecture \textup{(C)} is false. More precisely,
\[
 \boxed{\itf{4}{9}(132)=8>7=\itf{5}{10}(132).}
 \tag{2}
\]
The functions \(\itf{4}{9}\) and \(\itf{5}{10}\) are incomparable in pointwise
order, since also
\[
 \itf{4}{9}(36)=2<3=\itf{5}{10}(36).
 \tag{3}
\]
\end{theorem}

\begin{proof}
The defining recurrence gives the following exact iteration chains:
\begin{align*}
132&\xrightarrow{F_4}96\to70\to51\to37\to27\to20\to15\to11\to8,\\
132&\xrightarrow{F_5}100\to75\to56\to42\to32\to25\to19\to14\to10\to7.
\end{align*}
The Lean certificate checks every entry of the two recurrence tables through
index 132, proves their agreement with the infinite functions in (1), and
deduces (2) and (3). Since \(9>2\cdot4\), (2) contradicts (C).
The opposite strict inequalities at 132 and 36 prove incomparability.
Section~\ref{sec:formal} gives the exact formal statements.
\end{proof}

\section{An infinite family}

The defining recurrence gives a counterexample for every parameter in
the range below. The proof uses the inverse relation established in
the next lemma.
For \(k\ge1\), set
\[
 L_k(n)=n+\itf{k}{k-1}(n).
 \tag{4}
\]

\Needspace{10\baselineskip}
\begin{theorem}[Failure for every parameter \(k\ge4\)]
For each integer \(k\ge4\), let
\[
 p=k+3,\qquad B_k=\itl{k+1}{2k+2}(p).
\]
Then
\[
 \boxed{\itf{k}{2k+1}(B_k)>k+3
      =\itf{k+1}{2k+2}(B_k).}
 \tag{5}
\]
In particular the proposed pointwise ordering fails for every \(k\ge4\)
at the first integer \(j\) strictly greater than \(2k\).
\end{theorem}

\subsection{The inverse relation}

Equation (1) defines a unique function on all natural numbers with
\(0\le F_k(n)\le n\).
Indeed, after all values below \(n+1\) have been constructed, every argument in
\(\itf{k}{k}(n)\) remains at most \(n\).
The subtraction therefore gives a nonnegative value and preserves the bound.

\begin{lemma}
For \(k\ge1\), consecutive values of \(F_k\) differ by either zero or one.
Moreover
\[
 F_k(L_k(n))=n,\qquad F_k(L_k(n)+1)=n+1,
 \tag{6}
\]
and, for every \(r,n,m\in\N\),
\[
 \itf{k}{r}(n)\le m\quad\Longleftrightarrow\quad n\le\itl{k}{r}(m).
 \tag{7}
\]
In particular \(\itf{k}{r}(\itl{k}{r}(m))=m\).
\end{lemma}

\begin{proof}
The first difference is \(F_k(1)-F_k(0)=1\). Suppose the zero-or-one step property
has been proved up to input \(n+1\). Iterating \(F_k\) preserves two ordered
inputs at distance at most one within that range. The outputs remain bounded
by the inputs. Hence
\(\itf{k}{k}(n+1)-\itf{k}{k}(n)\in\{0,1\}\).
Subtracting the two instances of (1) proves the next step. This also proves
monotonicity, both for \(F_k\) and for all its iterates.

By (4), \(L_k(n+1)-L_k(n)\in\{1,2\}\). The first identity in (6) follows
by induction on \(n\). It holds at zero. If \(u=L_k(n)\) and \(F_k(u)=n\),
then (1) gives
\[
 F_k(u+1)=u+1-\itf{k}{k-1}(n)=n+1.
\]
Put \(v=L_k(n+1)\). If \(v=u+1\), the induction step is complete. If
\(v=u+2\), applying (1) once more gives
\[
 F_k(v)=v-\itf{k}{k-1}(n+1)=n+1.
\]
The displayed calculation of \(F_k(u+1)\) also gives the second identity in (6).
Monotonicity and (6) now imply
\(F_k(n)\le m\) exactly when \(n\le L_k(m)\).
Iteration proves (7). Iterating the first identity in (6) proves the final
claim. These arguments also apply at zero.
\end{proof}

\subsection{Recurrences for the inverse iterates}

For a fixed starting value \(p\), write \(a_t=\itl{k}{t}(p)\).
The cancellation identity in (6) and definition (4) imply, for all \(t\ge0\),
\[
 a_{t+k}=a_{t+k-1}+a_t.
 \tag{8}
\]
To see this directly, apply (4) to \(\itl{k}{t+k-1}(p)\) and cancel
\(k-1\) pairs of \(F_k\) and \(L_k\).

Now take \(k\ge4\), \(p=k+3\), and define
\[
 q_t=p+\frac{t(t+5)}2.
 \tag{9}
\]
The small values obtained from (1) are
\[
 F_k(n)=n-1\quad(2\le n\le k+2),\qquad F_k(k+3)=k+1.
 \tag{10}
\]
To prove these identities, first note that \(F_k(1)=1\).
For \(2\le n\le k+2\), the previous values decrease by one until
reaching the fixed point 1, so
\(\itf{k}{k}(n-1)=1\). At \(n=k+3\), the same argument gives
\(\itf{k}{k}(k+2)=2\). Induction therefore proves (10).

Using (10) and cancellation in (4), the initial segments of the two relevant
inverse orbits are
\[
 \begin{array}{ll}
 a_t=\itl{k}{t}(p)=q_t,&0\le t\le k-1,\\[2pt]
 b_t=\itl{k+1}{t}(p)=q_t,&0\le t\le k.
 \end{array}
 \tag{11}
\]
To calculate the increments, let $1\leq t\leq k-1$.
Cancellation of $t-1$ pairs gives
\[
 a_t-a_{t-1}=F_k^{\,k-1}(L_k^{\,t-1}(p))
 =F_k^{\,k-t}(p).
\]
The first application sends $p=k+3$ to $k+1$ by (10). Each of the remaining
$k-t-1$ applications subtracts one. All these inputs lie between $2$ and
$k+2$. Thus $a_t-a_{t-1}=k+1-(k-t-1)=t+2$.
For the other orbit, $1\leq t\leq k$ gives
\[
 b_t-b_{t-1}=F_{k+1}^{\,k}(L_{k+1}^{\,t-1}(p))
 =F_{k+1}^{\,k-t+1}(p).
\]
Now $p=k+3=(k+1)+2$ lies in the initial range of $F_{k+1}$.
Each of the $k-t+1$ applications subtracts one, so the increment is again
$k+3-(k-t+1)=t+2$. Both orbits start at $p$. Summing the increments
$3,4,\ldots,t+2$ proves (9) and (11).

Let \(S=\sum_{t=0}^{k-1}q_t\) and \(A_k=a_{2k+1}\). Summing (8) through one
block, then taking the next two steps, yields
\[
 A_k=3q_{k-1}+S+2q_0+q_1.
 \tag{12}
\]
In detail, \(a_{2k-1}=q_{k-1}+S\),
\(a_k=q_{k-1}+q_0\), and
\(a_{k+1}=q_{k-1}+q_0+q_1\). Adding the last two terms to the first gives (12).
The same calculation with recurrence order \(k+1\) gives
\[
 B_k=b_{2k+2}=3q_k+S+q_0.
 \tag{13}
\]
Consequently
\begin{align*}
 B_k-A_k
 &=3(q_k-q_{k-1})-q_0-q_1\\
 &=3(k+2)-(k+3)-(k+6)\\
 &=k-3>0.
 \tag{14}
\end{align*}
The sum \(S\) cancels, so its closed form is not needed.

\begin{proof}[Proof of Theorem 2]
By (6) iterated,
\(\itf{k+1}{2k+2}(B_k)=p\). Equation (14) says
\(B_k>\itl{k}{2k+1}(p)\). The equivalence (7) therefore gives
\(\itf{k}{2k+1}(B_k)>p\), proving (5).
\end{proof}


\section{Formal verification and reproducibility}\label{sec:formal}
Lean defines the functions on all natural numbers by well-founded recursion.
The theorem \path{validTable_correct} proves that a table satisfying the
recurrence agrees with the function throughout its index range.
The Lean kernel checks the tables for $k=4,5$ by reduction of \texttt{decide}.

The namespace \texttt{MathAndTheory.Hofstadter} contains the theorems
\path{pointwise_incomparable} and \path{lateIterateOrderConjecture_false}.
The foundation module proves the step property, both inverse identities,
and the iterate equivalence (7). The family module proves
\path{inverse_iterate_gap} and \path{general_counterexample}.
The latter requires only $k\geq4$. Every inverse identity it uses is proved
from the recurrence.

An independent arithmetic check computes the tables as follows.
Begin with $a_0=0$. At step $m\geq1$, set $u=m-1$.
Replace $u$ by $a_u$ exactly $k$ times, then append $m-u$.
All lookup indices are less than $m$ because $a_i\leq i$.
The calculation therefore uses only integers and terminates at each step.
This finite check supports the witness. Lean also proves the counterexample
for every parameter in Theorem 2.

\section{Relation to earlier work}
Conjecture (C) concerns more iterations than the proved base comparison
$F_k\leq F_{k+1}$ or the comparison at $j=k+1$.
Letouzey's preprint \cite{new}, deposited on 26 October 2026, proves the
earlier-index comparison and strict base ordering beyond an explicit bound.
Its Section 7.2 concerns $F_k^{k+1}\geq F_{k+1}^{k+2}$.

The associated Rocq development \cite{repo} contains a witness to
$F_k^j\not\geq F_{k+1}^{j+1}$ for $j\geq k+2$.
Its witness to failure of the opposite order applies when $j\leq2k$.
That range excludes $j=2k+1$, which is the index in Theorem 2.
The inverse relations used here are due to the earlier research,
although this formalization proves them again from the recurrence.
No claim of first priority is made.

The infinite family varies $k$. Theorem 2 gives one input for each parameter.
It does not assert infinitely many violating inputs for each fixed $k$
or classify all pairs of iteration indices. The finite incomparability
statement concerns the pair $F_4^9,F_5^{10}$.

\section*{AI assistance}
Dev D. Goyal directed the research and is responsible for publication.
ChatGPT (OpenAI), using GPT-6 with Pro reasoning, generated the mathematical
arguments and Lean formalizations. Codex assisted with formalization,
computations, and editing. This preprint has not been peer reviewed.

\bibliographystyle{plainnat}
\bibliography{references}
\end{document}
